\hbadness=9000
% Sprachunterstützung 
% Mehrere Sprachen für angepasstes Literaturverzeichnis (babelbib) 
\usepackage[english,german,ngerman]{babel}



% http://mrunix.de/forums/showthread.php?p=301707
% Abstand zwischen Absätzen in minipage-Umgebung:
\makeatletter
\newcommand*{\@minipagerestore}{\parskip@update}
\makeatother



%%%%%%%%%%%%%%%%%%%%%%%%%%%%%%%%%%%%%%%%%%%%%%%%%%%%%%%%%%%
%%% Verzeichnisse
\usepackage[
% tight,
]{shorttoc}
%
% % When you want to set only the table of contents of a document in two (or
% % more columns), there is one known way1: Add an
% % \addtocontents{toc}{\protect\begin{multicols}{2}}
% % before the \tableofcontents and an
% % \addtocontents{toc}{\protect\end{multicols}}
% % at the end of the document. This way your .toc will start with
% % \begin{multicols}{2} and end with \end{multicols}.
%
% \usepackage[toc]{multitoc}

% increase width for large page numbers in TOC
%% http://www.kronto.org/thesis/tips/format-toc.html
%% Because of the font change, the page number becomes too large for the
%% horizontal space LaTeX reserves for it by default. Without the following
%% redefines to fix it, this would cause the Chapter entry page numbers
%% to extend a few points into the right margin. The horror!

\makeatletter 
\renewcommand*{\@pnumwidth}{3.5em}
\renewcommand{\@tocrmarg}{4.5em}

\makeatother

% \usepackage{titletoc}
% \contentsuse{lomassnahmen}{lomassnahmen}
% \titlecontents{massnahmen}[0pt]{}{}{\titlerule}{}

% /part ohne Seitenzahl im TOC?
% http://mrunix.de/forums/showthread.php?t=61596
\makeatletter
\let\partbackup\l@part
\renewcommand*\l@part[2]{\partbackup{#1}{}}
\makeatother

\makeatletter
\let\chapterbackup\l@chapter
\renewcommand*\l@chapter[2]{\chapterbackup{#1}{}}
\makeatother


\usepackage{makeidx}
\setcounter{secnumdepth}{3}
% LaTeX-Begleiter S.676
\makeatletter
\renewenvironment{theindex}
   {
   \begin{WidePar}
   ~

   \begin{multicols}{2}[\chapter{\indexname}]
      \setlength\parindent{0pt}\setlength\parskip{0pt}\let\item\@idxitem
   }
   {
   \end{multicols}
   \end{WidePar}
   }
\makeatother

% Hat man den Wunsch, auf eine Fußnote mehrfach zu verweisen, bietet LATEX dafür keine direkte Möglichkeit.
% Man muss hier die folgende Definition, eingerahmt von \makeatletter und \makeatother, in den Vorspann des Dokuments einfügen.
% Im Text muss bei der Fußnote, auf die ein zweites mal verwiesen werden soll, ein Label (hier: \label{fn:wichtig})eingefügt werden.
% Mit \footref{fn:wichtig} verweist man dann ein weiteres mal auf sie.
% http://www2.informatik.hu-berlin.de/~piefel/LaTeX-PS/Archive-2004/V07-footnote.pdf
\makeatletter
\providecommand*{\footref}[1]{% \begingroup 
\unrestored@protected@xdef\@thefnmark{\ref{#1}}% \endgroup 
\@footnotemark}
\makeatother

\usepackage{prettyref}
\newrefformat{sec}{\emph{\ref{#1}}}
% \newrefformat{pkt}{\emph{Abs.\,\ref{#1}}}
% \newrefformat{lit}{\emph{Abs.\,\ref{#1}}}
\newrefformat{chp}{Kap.~\ref{#1}}
\newrefformat{fig}{\SignFigureSmall~\ref{#1}}
\newrefformat{tab}{\SignTableSmall~\ref{#1}}
% \newrefformat{tab}{{\sffamily\mdseries\color{DefaultB}{\scalefont{0.7}\ovalbox{T}}}\,\ref{#1}}
% \newrefformat{m}{{\sffamily\mdseries\color{DefaultB}{\scalefont{0.7}\ovalbox{M}}}\,\ref{#1}}
% \newrefformat{m}{\SignMassnahmeSmall\,\ref{#1}}
\newrefformat{m}{\ref{#1}}
\newrefformat{AASS-topic}{{\AASS}~\ref{#1}}
\newrefformat{topic}{\Pointinghand\,\ref{#1}}

\usepackage[AASS,stylelayout]{aass-common}
\usepackage{aass-units}
\usepackage{aass-fonts}
\usepackage{aass-features}
\usepackage{aass-text}
\usepackage{aass-layout}










%%%%%%%%%%%%%%%%%%%%%%%%%%%%%%%%%%%%%%%%%%%%%%%%%%%%%%%%%%%
%%% Seitenlayout
\usepackage[
verbose,
a5paper,
asymmetric,
% reversemp,
includemp,   % incl. marginpar
includehead, % incl. header
includefoot, % incl. footer
inner  = 30pt,
outer  = 30pt,
top    = 7mm,
bottom =  7mm,
bindingoffset=22.50787pt,%0.50787+42.67912+2 → Rest=552pt 
marginparsep=0pt,
marginparwidth = 0pt,
% marginparpush=0pt,
footskip = 7mm,
% % headheight=0.5cm,
% % headsep=0.5cm,	
% showframe,
]
{geometry}


\setlength\columnsep{15pt}%7mm
\setlength\columnseprule{0.4pt}
\setlength\premulticols{10\baselineskip}
\renewcommand{\columnseprulecolor}{\color{DefaultGrey}}

%%% Tables

\usepackage{boxedminipage}

\usepackage{layout}

\newlength{\WideParWidth}
\setlength{\WideParWidth}{\textwidth+\marginparwidth+\marginparsep}

\newlength{\ThirdGrossWidth}
\setlength{\ThirdGrossWidth}{(\WideParWidth)/3}

\newlength{\ThirdNetWidth}
\setlength{\ThirdNetWidth}{(\WideParWidth-(\columnsep*2))/3}

\newlength{\Textspalte}
\setlength{\Textspalte}{\textwidth}

\newlength{\Randspalte}
\setlength{\Randspalte}{\marginparwidth}

\newlength{\Trenner}
\setlength{\Trenner}{\marginparsep}

\newlength{\Korrektur}
\setlength{\Korrektur}{{\fboxrule} + {\fboxsep}}
% \setlength{\Korrektur}{0pt}

% \usepackage{scrindex}
\makeindex













%%%%%%%%%%%%%%%%%%%%%%%%%%%%%%%%%%%%%%%%%%%%%%%%%%%%%%%%%%%
%%% Pargragraphensatzregeln
%Dann darst du die parskip und \parindent gar nicht verändern, sondern Koma Script in der Klasse die entsprechende Option übergeben. 
%\setlength{\parindent}{0pt}
%\setlength{\parskip}{5pt}
\raggedbottom


% \@beginparpenalty - Strafpunkte für einen Umbruch vor einer Liste 
% \@endparpenalty - Strafpunkte für einen Umbruch nach einer Liste 
% \@itempenalty - Strafpunkte für einen Umbruch zwischen Listeneinträgen 
\makeatletter% --> Wiki 
\@itempenalty=500 
\makeatother% --> Wiki


% % (Anm.: Werte von -10000 bzw. 10000 gelten bei TeX als
% % negativ unendlich bzw. positiv unendlich.) 
% % Textsatz: interner Parameter zur Steuerung des Zeilenumbruchs. Er
% % enthält die maximale Minuspunktzahl, die ein Absatz beim
% % Umbruch haben darf, wobei diese für den ersten Durchgang
% % ohne Trennungen gilt. Setzt man \pretolerance=10000, gilt
% % der Absatz immer als akzeptiert; es wird also nicht
% % getrennt.
% \pretolerance=100
% % Bewertungsgrenzwert für noch zu akzeptierende schlecht
% % umgebrochene Zeilen, und zwar mit erlaubten Trennungen
% % (zweiter Durchgang). Wird \tolerance vergrößert, so kommt es
% % auch zu mehr Leerraum zwichen einzelnen Wörtern. Bei
% % \tolerance=10000 wird auf den dritten Durchgang verzichtet.
% \tolerance=300
% % Grenze, ab der eine overfull hbox gemeldet wird. (Wert, den
% % eine Zeile mit der Länge \hsize überschreiten darf.)
% \hfuzz=0.1pt
% % Grenzwert, ab dem die Überfüllung einer \vbox protokolliert
% % wird.
% \vfuzz=0.1pt
% % Grenzwert für »schlechte« Zeilen, bzw. Boxen, der angibt,
% % ab welcher Negativbewertung eine \hbox protokolliert wird.
\hbadness=9999
% % Grenzwert, ab dem eine »schlechte« \vbox protokolliert wird.
\vbadness=9999
% % Definiert zusätzlichen dynamischen Leerraum, der falls er
% % von 0 Punkten verschieden ist, beim Absatzumbruch einen
% % dritten Versuch bewirkt, wenn bei den beiden ersten
% % Durchläufen kein Umbruch ohne eine overfull hbox gelungen
% % ist. \emergencystretch definiert dann den Leerraum, der
% % innerhalb einer Zeile zusätzlich verteilt werden darf.
% % \emergencystretch=20pt
\emergencystretch=5em

% Strafpunkte (penalties)
% 
% Zusätzlich zu den Dehnungspunkten gibt es noch sogenannte
% Strafpunkteregister (penalties). Sie enthalten
% Minuspunkte, die für bestimmte Umbruchkonstellationen zu
% dem »Unschönheitsfaktor« hinzuaddiert werden und damit
% letztendlich die Bewertung dieser Konstellation bilden: Es
% wird diejenige Alternative gewählt, die die wenigsten
% Strafpunkte ergibt. 
% 
% % Grundlast, die für jede Zeile
% % berechnet wird. Wird dieser Wert erhöht, so versucht das
% % Programm, die Zeilenanzahl für einen Absatz möglichst
% % klein zu halten.
\linepenalty=10
% % Wird für eine einfache Trennung beim Absatzumbruch aufgerechnet.
\hyphenpenalty=40
% % Wird für einen Zeilenwechsel nach einem expliziten
% % Trennstrich aufgerechnet.
% \exhyphenpenalty=50
% % Mathematiksatz: Gibt es für das Trennen einer Formel im
% % text-style nach einem binären Operator.
% \binoppenalty=700
% % Wird beim Trennen einer mathematischen Formel im
% % text-style (nach einer Relation) aufgerechnet.
% \relpenalty=500
% % Wird beim Seitenumbruch vergeben, falls die erste Zeile
% % eines Absatzes allein auf der vorangehenden Seite verbleibt.
% % Disable single lines at the start of a paragraph (Schusterjungen)
\clubpenalty=150
% % Wird beim Seitenumbruch aufgerechnet, falls die letzte Zeile
% % gerade noch auf die nächste Seite umgebrochen wird.
% % Disable single lines at the end of a paragraph (Hurenkinder)
\widowpenalty=150
% % Wird während des Seitenumbruchs für einen Seitenwechsel
% % angerechnet, falls die letzte Zeile eines Absatzes gerade
% % noch auf die nächste Seite kommt, und zwar in dem
% % Spezialfall, dass direkt danach eine hervorgehobene Formel
% % folgt.
% \displaywidowpenalty=50
% % Wird während des Seitenumbruchs für einen Seitenwechsel
% % angerechnet, bei dem die letzte Zeile einer Seite mit einer
% % Trennung endet.
\brokenpenalty=40
% % Für zwei aufeinanderfolgende Zeilen, wobei in beiden getrennt wird.
\doublehyphendemerits=100
% % Wird für die Trennung in der vorletzten Zeile eines Absatzes aufgerechnet.
\finalhyphendemerits=40
% % Wird beim Absatzumbruch für zwei »visuell inkompatible«
% % Zeilen aufgerechnet. Dies sind zwei aufeinanderfolgende
% % Zeilen, bei denen die eine mit extra viel Leerraum
% % durchgeschossen wurde und die andere mit wenig.
\adjdemerits=9000
% % Wird für das Umbrechen einer mit \displaylines
% % entstandenen Formelfolge über eine Seitengrenze hinweg
% % berechnet.
% \interdisplaylinepenalty=100
% % Wird für das Umbrechen einer Fußnote über Seitengrenzen
% % hinweg berechnet. Der Befehl \footnote setzt die
% % \interlinepenalty auf den angegebenen Wert.
% \interfootnotelinepenalty=100
% % Für den Umbruch eines Absatzes über eine Seitengrenze.
\interlinepenalty=50
% % Für einen Seitenwechsel direkt vor einer Formel im display-style.
% \predisplaypenalty=10000
% % Für einen Seitenwechsel direkt nach einer Formel im display-style.
% \postdisplaypenalty=0


\usepackage{fancybox}


% User page numbering with chapter number in it

\makeatletter
\newcommand*{\greek}[1]{%
  \expandafter\@greek\csname c@#1\endcsname
}
\newcommand*{\@greek}[1]{%
$\ifcase#1\or\alpha\or\beta\or\gamma\or\delta\or\varepsilon\or\zeta\or\eta\or\theta\or\iota\or\kappa\or\lambda\or\mu\or\nu\or\xi\or o\or\pi\or\varrho\or\sigma\or\tau\or\upsilon\or\phi\or\chi\or\psi\or\omega\else\@ctrerr\fi$%
}
\makeatother

\renewcommand{\thechapter}{\greek{chapter}}

\renewcommand{\thechapter}{\Alph{chapter}}

% \renewcommand{\thechapter}{{\color{DefaultB}MK11.1/}\arabic{chapter}}
% \renewcommand{\thechapter}{{\color{DefaultB}MK11.1/}\arabic{chapter}}

% \renewcommand{\thechapter}{{\ttfamily\Alph{chapter}}}
\renewcommand{\thesection}{{\color{Indigo}{\SignMassnahme}\,}{\color{Indigo}\thechapter.\arabic{section}}}
% \renewcommand{\thesubsection}{{\color{DefaultA}\thesection .\arabic{subsection}}}
% \renewcommand{\thesubsection}{{\color{aassgrey}\thesection.}\arabic{subsection}}
% \renewcommand{\thesubsubsection}{{\thesubsection .\alph{subsubsection}}}



% No bottom titles !!!
\usepackage[nobottomtitles]{titlesec}
\renewcommand{\bottomtitlespace}{.10\textheight}

% \titleformat{\section}[rightmargin]{\usefont{T1}{lmss}{bx}{n}\selectfont\LARGE\color{DefaultA}}{}{0em}{}
% \titleformat{\section}{\large\scshape\raggedright}{}{1em}{}


%%%%%%%%%%%%%%%%%%%%%%%%%%%%%%%%%%%%%%%%%%%%%%%%%%%%%%%%%%%
%%% Floats
\usepackage{wrapfig,picins}

% watches margin paragraphs and
% gives warnings
%\usepackage{mparhack} 




\usepackage{hyperref}

% \newcommand{\url}[1]{\ttfamily #1}


% \usepackage[grey,helvetica]{quotchap}
% %\renewcommand{\sectfont}{\large\color{DefaultA}\sffamily\fontfamily{phv}\selectfont}
% \definecolor{chaptergrey}{rgb}{0.788235294,0,0.08627451}
% \definecolor{chaptergrey}{named}{DefaultContrast}
% % \definecolor{chaptergrey}{rgb}{0.788235294,0,0.08627451}
% % \newenvironment{savequote}{}{}
% 
% % \usepackage[Rejne]{fncychap}
% % \ChNameVar{\centering\LARGE\sf} 
% % \ChNumVar{\Huge\color{DefaultB}} 
% % \ChTitleVar{\centering\Large\sf\color{DefaultA}}
% % \ChRuleWidth{2pt} 
% % \ChNameUpperCase
% % % \ChNumVar{\fontsize{76}{80}\usefont{OT1}{lmdh}{m}{n}\selectfont\color{DefaultB}}
\addtokomafont{chapter}{\Huge\sffamily\bfseries\color{DefaultA}}

% \addtokomafont{chapter}{\fontfamily{lmdh}\selectfont\Huge\color{DefaultA}}









\usepackage{minitoc} 
%The quotchap package should be loaded BEFORE the minitoc package.
\newcounter{DefaultMinitocDepth}
\setcounter{DefaultMinitocDepth}{3}
\mtcsetdepth{minitoc}{\value{DefaultMinitocDepth}}

\nomtcrule

% \setlength{\mtcindent}{0pt}
% \setlength{\stcindent}{0pt}
% \mtcindent{0pt}
% \stcindent{0pt}
\renewcommand{\mtcfont}{\sffamily\small}
\renewcommand{\mtcSfont}{\sffamily\small\bfseries}
\renewcommand{\mtcSSfont}{\sffamily\scriptsize}
\renewcommand{\mtcSSSfont}{\sffamily\scriptsize}
% \mtcsetoffset{minitoc}{-4em}
% \mtcsetformat{mini-table}{mtcindent}{0pt}

\setlength{\mtcindent}{0pt}
\mtcsetfeature{minitoc}{open}{\kern1sp\vspace*{-.1ex}\begin{multicols}{2}[\kern-2.5ex]}
\mtcsetfeature{minitoc}{close}{\end{multicols}\kern-2.ex}
% \mtcsetfeature{minitoc}{open}{\begin{multicols}{2}}
% \mtcsetfeature{minitoc}{close}{\end{multicols}}
\mtcsetfeature{minitoc}{before}{\begin{WidePar}\raggedcolumns}
\mtcsetfeature{minitoc}{after}{\flushcolumns\end{WidePar}}
\mtcsetoffset{minitoc}{-0.7em}
\setlength{\mtcindent}{-0.7em}
% \mtcsetformat{minitoc}{tocrightmargin}{2cm}

% % multitoc is incompatible with float package
% \usepackage[toc]{multitoc}
% \renewcommand{\multicolumntoc}{2}







% \usepackage[noauto]{chappg} % nach hyperref!
% 
% \pagenumbering{bychapter}
% \renewcommand{\chappgsep}{/}


%%%%%%%%%%%%%%%%%%%%%%%%%%%%%%%%%%%%%%%%%%%%%%%%%%%%%%%%%%%
%%% Beschriftungen
% \DeclareRobustCommand{\SignTable}{{\noindent\sffamily\mdseries\color{DefaultB}\ovalbox{T}}}
% \DeclareRobustCommand{\SignTable}{{\noindent\scalefont{1.5}$\boxbar$}}
\DeclareRobustCommand{\SignTable}{{Tab.}}
% \DeclareRobustCommand{\SignTableSmall}{{\scalefont{0.7}\SignTable}}
\DeclareRobustCommand{\SignTableSmall}{{\SignTable}}

\DeclareRobustCommand{\SignFazit}{\noindent{\fontfamily{lmr}\selectfont >>}}
% \DeclareRobustCommand{\SignFazit}{\noindent{\color{DefaultContrastA}\bfseries{\VarFlag ~}}}
\DeclareRobustCommand{\SignFazitSmall}{\SignFazit}

% \DeclareRobustCommand{\SignLeitsatz}{{\noindent\sffamily\mdseries\color{DefaultB}\ovalbox{L}}}
\DeclareRobustCommand{\SignLeitsatz}{\SignFazit}
\DeclareRobustCommand{\SignLeitsatzSmall}{\SignFazit}

% \DeclareRobustCommand{\SignFigure}{{\noindent\sffamily\mdseries\color{DefaultB}\ovalbox{A}}}
% \DeclareRobustCommand{\SignFigureSmall}{{\scalefont{0.7}\SignFigure}}
% \DeclareRobustCommand{\SignFigure}{{\noindent\color{DefaultB}\BigSquare\llap{{\VarIceMountain}}}}
\DeclareRobustCommand{\SignFigure}{{Abb.}}
\DeclareRobustCommand{\SignFigureSmall}{{\SignFigure}}

% \DeclareRobustCommand{\SignMassnahme}{{\noindent\sffamily\mdseries\color{DefaultB}\ovalbox{M}}}
% \DeclareRobustCommand{\SignMassnahmeSmall}{{\scalefont{0.7}\SignMassnahme}}

\DeclareRobustCommand{\SignMassnahme}{{\noindent\color{Indigo}\dsmedical}}
\DeclareRobustCommand{\SignMassnahmeSmall}{{\SignMassnahme}}

% \DeclareRobustCommand{\SignMassnahme}{{\noindent\color{DefaultB}\BigSquare\llap{{\sffamily\scalefont{0.9}M}}}}
% \DeclareRobustCommand{\SignMassnahmeSmall}{{\SignMassnahme}}








%%%%%%%%%%%%%%%%%%%%%%%%%%%%%%%%%%%%%%%%%%%%%%%%%%%%%%%%%%%
%%% Bibliographie


\newcommand{\ChapterBibliography}{\relax}


\usepackage[version=3]{mhchem}
\mhchemoptions{textfontcommand=\sffamily}
\mhchemoptions{mathfontcommand=\mathsf}




% various packages for convenience


\usepackage{rotating}
\usepackage{subfig}
\usepackage{needspace}


% Test PLAYGROUND

\newcommand{\FullRef}[1]{{\ref{#1} (S.\,\pageref{#1})}}
\newcommand{\FullPrettyRef}[1]{{\prettyref{#1}, Seite \pageref{#1}}}

\newcommand{\TabItemI}[1]{\begin{inparaitem}\item #1 \end{inparaitem}}
\newcommand{\TabItemII}[1]{\begin{inparaitem}\begin{inparaitem}\item#1\end{inparaitem}\end{inparaitem}}
\newcommand{\TabItemIII}[1]{\begin{inparaitem}\begin{inparaitem}\begin{inparaitem}\item#1\end{inparaitem}\end{inparaitem}\end{inparaitem}}



% \setkomafont{footnote}{\sffamily}
\deffootnotemark{\textsuperscript{\sffamily\thefootnotemark}}
% \setlength{\footnotesep}{1pt}
\usepackage[hang,multiple]{footmisc}
% \renewcommand\footnotelayout{\sffamily}



\newcommand{\papertexinexpandedtitle}[1]{
\begin{minipage}{.95\textwidth}
\begin{center}
\noindent\Large\textbf{#1}
\end{center}
\end{minipage}
}

\newcommand{\expandedtitle}[2]{
% \end{multicols}
% \begin{center}
\setlength{\fboxsep}{5pt}
\setlength{\shadowsize}{2pt}
\ifthenelse{\equal{#1}{shadowbox}}{%
\shadowbox{%
\papertexinexpandedtitle{#2}%
}%
}{}
\ifthenelse{\equal{#1}{doublebox}}{%
\doublebox{%
\papertexinexpandedtitle{#2}%
}%
}{}
\ifthenelse{\equal{#1}{ovalbox}}{%
\ovalbox{%
\papertexinexpandedtitle{#2}%
}%
}{}
\ifthenelse{\equal{#1}{Ovalbox}}{%
\Ovalbox{%
\papertexinexpandedtitle{#2}%
}%
}{}
\ifthenelse{\equal{#1}{lines}}{
{\color{DefaultB}\hrule}
\vspace*{1em}
\begin{center}
\noindent\sffamily\textit{#2}
\end{center}
\vspace*{1em}
{\color{DefaultB}\hrule}
}{}
\ifthenelse{\equal{#1}{roundbox}}{
% {\color{DefaultB}\hrule}
% \vspace*{8pt}
% \begin{center}
% \\[0.5\baselineskip]\BoxRoundedBorderFilled{DefaultD}{White}{\center\hspace{0pt}\normalsize{\SlideHeading{#2}}}\\[0.5\baselineskip]
\BoxRoundedBorderFilled{DefaultD}{White}{\center\hspace{0pt}\bfseries#2}
% \vspace{-0.5\baselineskip}
% \end{center}
% \vspace*{8pt}
%{\color{DefaultB}\hrule}
}{}
\ifthenelse{\equal{#1}{simple}}{
\begin{center}
\bfseries#2
\end{center}
}{}
% \end{center}
% \begin{multicols}{\papertex@ncolumns{}}
% \ifnum \papertex@ncolumns > \minraggedcols
% \raggedFormat
% \fi
}

\renewcommand{\FontTableBase}{\rmfamily}
\renewcommand\TabHeading[1]{{\color{DefaultB}\textbf{#1}}}
\renewcommand\TabSub[1]{{\bfseries #1}}


\usetikzlibrary{shapes,calc}%Tbl of elements



% Randnummernzähler einrichten 
\newcounter{randnr} 
\newcommand{\newrandnummer}{\refstepcounter{randnr}} 
\newcommand{\rdnr}{\newrandnummer\marginnote{{\sffamily\bfseries\scriptsize\therandnr}}} 

\makepagenote

% \newcommand{\MultiColMassnahmenStart}[1]{\addtocontents{lomassnahmen}{\protect\begin{multicols}{#1}}}
% \newcommand{\MultiColMassnahmenEnd}{\addtocontents{lomassnahmen}{\protect\end{multicols}}}
% \newcommand{\MultiColFazitStart}[1]{\addtocontents{lofazit}{\protect\begin{multicols}{#1}}}
% \newcommand{\MultiColFazitEnd}{\addtocontents{lofazit}{\protect\end{multicols}}}
  
\newcommand{\MultiColListStart}[2]{\addtocontents{#1}{\protect\begin{multicols}{#2}}}
\newcommand{\MultiColListEnd}[1]{\addtocontents{#1}{\protect\end{multicols}}}
\newcommand{\Sign}[1]{{\sffamily#1}}


\newcommand{\KAMEL}[1]{{\color{DefaultB}\small\textcircled{\ttfamily\bfseries\scriptsize#1}}}
\newcommand{\TxKAMEL}{(S-)KAMEL}
\newcommand{\SpeechMargin}[1]{\marginpar{{\small\Speech{#1}}}}
\newcommand{\SpeechTextMargin}[1]{\marginpar{{\small\Speech{#1}}}\Speech{#1}}
\newcommand{\Study}[1]{\textit{#1}}
\newcommand{\FontTypewriter}{\fontfamily{lmtt}\fontseries{lc}\selectfont}








\newcommand{\VARIANT}[2]{{\opt{STATUS-DRAFT}{\color{SteelBlue}\sffamily\textbf{///~\uline{#1}$\rightarrow$}}\opt{#1}{#2}\opt{STATUS-DRAFT}{$\leftarrow$\textbf{\uline{#1}~///}}}}
\usepackage[autostyle,german=guillemets]{csquotes}